\section{Introduction}
Language models (LMs) that can predict the next word for input text are a powerful tool for many applications. In mobile keyboard applications, LMs are deployed on device to support various features (e.g., auto correction, smart completion and suggestion, and next word prediction) to improve users’ typing experience. On-device LMs are typically small (less than ten million parameters) due to latency requirement and limited on-device resources. Their performance can be significantly improved by training from user data~\citep{hard2018gboard,xu2023federated}; recent work~\citep{wang2023can,wu2024prompt} shows the necessity of training on user data to achieve high utility even when we can access large-scale web data and pre-trained large LMs with billions of parameters. 

As mobile-keyboard user data can be highly privacy sensitive, differential privacy (DP)~\citep{dwork2006calibrating,dwork2014algorithmic} and federated learning (FL)~\citep{mcmahan2017fedavg,kairouz2019advances} have emerged as best practices for such models. 
DP provides a mathematical formulation to upper-bound the memorization of an individual’s information in model training. FL minimizes data exposure by aggregating focused model updates from decentralized data stored only on user devices. DP and FL are combined when training on-device language models in production mobile keyboard applications\ifthenelse{\boolean{acl}}{~\citep{xu2023federated}}{~\citep{xu2023federated,choquette2023amplified,wu2024prompt}}. Applying DP in a production cross-device FL system is challenging as many DP algorithms require specific pattern of sampling training data to achieve strong privacy-utility trade-off. However, a cross-device FL system has limited control of sampling as clients can only participate in training when local criteria (e.g., charging, idle, and connected to an unmetered network) are satisfied~\citep{bonawitz2019towards,huba2022papaya}. 
Recently, DP-Follow-the-Regularized-Leader (DP-FTRL) algorithms~\citep{kairouz21practical,choquette2023amplified} have achieved superior privacy-utility trade-off with simpler client participation requirements, and are used in practice in FL systems~\citep{xu2023federated,zhang2023private}. 

Instead of requiring uniform or Poisson sampling of devices as in previous work~\citep{abadi2016deep,mcmahan2017learning}, DP-FTRL uses minimum separation (min-sep) to characterizes the participation pattern\ifthenelse{\boolean{acl}}{}{~\citep{kairouz21practical,xu2023federated}}. Min-sep is the smallest number of rounds between the consecutive participation of a client, and smaller min-sep necessitates adding more noise to achieve a desired DP guarantee. Min-sep is enforced in the FL system by implementing a timer on each device so that a device only becomes eligible for training if a certain period of time (e.g., three days) has passed since their last participation.
DP-FTRL algorithms leverage correlated noise mechanisms such as tree aggregation (\treeagg)~\citep{kairouz21practical} or matrix factorization (MF)~\citep{choquette2023amplified} with the client participation pattern in FL. The banded MF (\bandmf) mechanism\ifthenelse{\boolean{acl}}{}{~\citep{choquette2023amplified}} pre-computes matrices to generate  correlated noise from independent noise to achieve stronger DP guarantees than the \treeagg mechanism\ifthenelse{\boolean{acl}}{}{~\citep{kairouz21practical}}. \bandmf is superior when the number of rounds and min-sep can be (accurately) estimated before training to optimize matrices. However, min-sep is only known after training with time-based separation as many system factors may potentially affect training time~\footnote{Despite being more complicated in FL systems, it is possible to enforce round-based separation so that a device only becomes eligible for training if min-sep rounds has passed since their last participation. However, it is still challenging to pre-specify min-sep before training due to the dynamics of client availability and population size. If the target min-sep is too large, training might halt because of lacking eligible devices. If the target min-sep is too small, the MF mechanism is not optimal for the correlated noise generation.}.  Furthermore, \bandmf consumes more memory for noise generation, and hence is used less often than \treeagg in practice. 

In this work, we focus on the challenges of achieving strong DP guarantees in training production LMs in a cross-device FL system. We discuss how we extend the recent theoretical advancements of the Buffered Linear Toeplitz (\blt) mechanism from single participation~\citep{mcmahan2024efficient} to multi-participation scenarios, and adapt \blt to DP-FTRL. We apply \blt-DP-FTRL to FL in practice, demonstrating its advantages in flexibility, ease of use, and privacy-utility trade-offs. The \blt-DP-FTRL algorithm offers flexibility in handling varying numbers of training rounds, and robustness to a wide range of min separation between user participations. Furthermore, \blt-DP-FTRL simplifies the correlated noise generation process and reduces memory by exploiting the parameterization of the Toeplitz matrices. 
We empirically evaluate \blt-DP-FTRL on the StackOverflow benchmark dataset and across four on-device LM training tasks in a production FL system. Our \blt-DP-FTRL achieves better privacy-utility trade-off compared to the widely used \treeagg mechanism\ifthenelse{\boolean{acl}}{}{~\citep{kairouz21practical}}, and comparable results compared to the state-of-the-art \bandmf mechanism\ifthenelse{\boolean{acl}}{}{~\citep{choquette2023amplified}}. 
\blt-DP-FTRL exhibits desirable robustness properties in practice, offering a practical and effective solution for achieving strong DP in real-world FL systems.
